% Part: set-theory
% Chapter: ordinals
% Section: milestone

\documentclass[../../../include/open-logic-section]{subfiles}

\begin{document}
\olfileid{sth}{ordinals}{zfm}
\olsection{$\ZFminus$: একটি মাইলফলক}

প্রতিস্থাপনকে কীভাবে ন্যায্যতা দেওয়া যায়, আদৌ যায় কি না,
সে প্রশ্নের উত্তর সহজ নয়। তাই আলোচনা
\olref[replacement][]{chap}-এর জন্য রেখে দিচ্ছি। তবে
প্রতিস্থাপন যোগ করার সঙ্গে সঙ্গে আমরা আরেকটি গুরুত্বপূর্ণ
মাইলফলকে পৌঁছেছি: এখন $\ZFminus$ তত্ত্বের জন্য প্রয়োজনীয়
সব স্বতঃসিদ্ধ আমাদের হাতে আছে। বিস্তারিতভাবে:

\begin{defn}
$\ZFminus$ তত্ত্বের স্বতঃসিদ্ধগুলি হলো: সদস্যভিত্তিক সমতা, সংযোগ,
যুগল, ঘাত সেট, অসীমতা এবং পৃথকীকরণ ও প্রতিস্থাপন
ছকের সব দৃষ্টান্ত। অন্যভাবে, $\Zminus$-এর সঙ্গে
প্রতিস্থাপন যোগ করলেই $\ZFminus$ পাওয়া যায়।
\end{defn}
\noindent
এটি \emph{জার্মেলো--ফ্রেঙ্কেল} সেটতত্ত্বের নাম
(\emph{একটি বিষয় বাদ দিয়ে}; সেটি পরে আসবে)।
ফ্রেঙ্কেলের নামটি যুক্ত হয়েছে, কারণ \citeyear{Fraenkel1922}-এ
প্রতিস্থাপনের প্রণয়ন তাঁর নামে স্বীকৃত; যদিও প্রথম
সুনির্দিষ্ট প্রণয়নটি \citet{Skolem1922}-এর।

\end{document}
